% Abstract.

\begin{abstract}
Language-model agents are increasingly used as persistent \emph{coworkers} that assist users across multiple working days. During such workflows, the surrounding environment may change independently of the agent: new emails arrive, calendar entries shift, knowledge-base records are updated, and evidence appears across images, scanned PDFs, audio, video, and spreadsheets. Existing benchmarks do not adequately evaluate this setting because they typically run within a single static episode and remain largely text-centric. We introduce \bench{}, a benchmark for coworker agents built around multi-turn multi-day tasks, a stateful sandboxed service environment whose state evolves between turns, and rule-based verification. The current release contains 100 tasks across 13 professional scenarios, executed against five stateful sandboxed services (filesystem, email, calendar, knowledge base, spreadsheet) and scored by 1{,}537 deterministic Python checkers over post-execution service state; no LLM-as-judge is invoked during scoring. We benchmark seven frontier agent systems. The strongest model reaches 75.8 weighted score, but the best strict Task Success is only 20.0\%, indicating that partial progress is common while complete end-to-end workflow completion remains rare. Turn-level analysis shows that performance drops after the first exogenous environment update, highlighting adaptation to changing state as a key open challenge. We release the benchmark, evaluation harness, and construction pipeline to support reproducible coworker-agent evaluation.
\end{abstract}
